\beginlecture{24}{17.07.2025}

\chapter{Review}

\section{General information}

\begin{itemize}
    \item 30-35 minutes
    \item he asks questions, and we answer
    \item more about general ideas and concepts
    \item write down statements, proof ideas (not 30 min technical details)
    \item Connections, what do we remember?
    \item We can write down some stuff, but not too much 
    \item Variable focus, but cover most stuff per exam
    \item Each exam different (also order) a bit of a random process. Last 2 week more like bonus 
    \item Cover every / most topic theme a bit at least
\end{itemize}

\section{Introduction to optimal control}
\begin{itemize}
    \item State space models and dynamics 
    \item Discrete first and continuous view
    \item Linear state space model examples connects to nicer results for linear function approximation
    \item ODE view,  good start stuff until here 
    \item Cost / reward functions  
    \item Stability and equilibria, Lyapunov methods, how do we understand stability (different versions) 
    \item Lyapunov functions (inf compactness (boundedness, control over sublevelsets)), other stuff a bit more interesting
    \item Results from Lyapunov functions for stability, convergence 
    \item Comparative theorem: connection from Lyapunov function to cost functions implies basis of everything (some proofs were nice!)
    \item \begin{itemize}
        \item started with fixed policy 
        \item finding policies (DP / Bellman equations)
        \item principal of optimality, helps reasonable assumption?  
    \end{itemize}
    \item Function representation, state action view (Q function), more flexibility for the algorithm
    \item If I have value function, I need a more specific model, whereas Q functions are more model agnostic? RL view, does not know model, control system often assumes it
    \item What is reinforcement learning hints 
    \item value and policy iteration, motivation by DP and proofs. In some form everything is either value or policy iteration 
    \item infinite Horizon (discount), and other formulations, model predictive control (more a bonus)
    \item continuous view of Lyapunov more intuitive (drift condition \(\frac{d}{dt} V(x)<0\), \(\langle \nabla V(x),f(x)\rangle<0\)), proof was inf compact 
    \item repeated optimal control in continuous time (to bellman equation? Nice to know) 
\end{itemize}

\section{ODE for algorithm design}
\begin{itemize}
    \item Root finding, describe / construct algorithm by some vector field 
    \item Connection between discretization, gradient descent, ODE scenarios
    \item ODE theory: Picard Lindelöf, Grönwall lemma
    \item Lemmas providing bounds, used later! Used either continuously or via Euler discretization 
    \item Convexity aspects, continuous differentiable + coercive to get convergence 
    \item Global asymptotic stability results and ODE view of gradient descent 
    \item PL-inequality condition and lower bound on gradient to get limit point behavior
    \item L-smooth + PL to get convergence and more specific statements on step size for Euler discretization
    \item L-smooth inequality and PL are bounds from two sides! 
    \item Properties of loss function to stability and convergence 
    \item What does all of this have to do with data? Quasi stochastic view (we later saw real stochastic view)
    \item Probing signal (without being too specific, example: sin and saw, similar to MC)
    \item  \(f,\bar{f}\) and their connection
    \item Comparing different views (thetas?)
    \item Quasi stochastic approximation ODE: one of the interesting ideas was the time scaling, this is a building block to compare views
    \item Same ODE but different timescale, properties: key thing was a pair of identities
    \item QSA1-QSA3, A going to zero but not to quickly, Lipschitz continuous, \(\bar{f}\) and \(f\) are reasonably connected 
    \item With that show bounds at the end: boundedness implies convergence  
    \item Question: how do get boundedness? Lyapunov comes to rescue with ultimate boundedness
    \item To get to that: QSV assumption
    \item This connected the ODE algorithm design to the earlier stuff
\end{itemize}
Now we have theoretical foundations!


\section{Value and \(Q\)-function approximation}
\begin{itemize}
    \item How can we get from trajectories and data to first algorithms?
    \item Bellman error reinterpret as temporal difference error
    \item TD as intermediate measurement, usually worked with \(Q\) function
    \item Mean squared bellman error 
    \item Algorithm: based on linearization solve LSE
    \item We saw several reasons for exploration: View that long trajectory after some point TD gives not much information (goes to zero), LSTD with restarts to get broader coverage 
    \item First LSTD was formulated in full approximation space, further view: projected bellman equation 
    \item Instead of bellman operator, introduced projected bellman operator: $\hat{h}=\hat{T}(\hat{h})=P_\cH(T(\hat{h}))$ or \(0=P_G(\hat{h}-T(\hat{h}))\)
    \item Under linearity assumptions and $\cH=G$ both views agree
    \item Started with linear approximation results
    \item Galerkin relaxation
    \item Going a bit back .. If we solve in the linear approximation with gradient methods, we can write Bellman error via gradient of loss function. Replace gradient by eligibility vectors    
    \item Formulate TD\(\lambda\), eligibility vectors and \(\lambda\), how do I take into account how temporal difference plays a role? \(\lambda=1,0\) gives to extremes.
    \item We later saw TD\((\lambda)\) has to do with bias variance, at this stage only as basic RL algorithm formulated in ODE setting 
    \item TD\((\lambda)\) in linear approximation scheme, can generalize, but can't write down ODE anymore 
    \item Derived a view algorithm like networks, batched versions ... 
    \item There we could alt least in the linear function approximation we can draw some connections, in non-linear versions, where we can't get much results 
    \item TD schemes: saw different ways in the policy evaluation? Galerkin as generalization, but breaking down the theory? 
    \item Additional observations: Observations and ODE approximation. Looking at specific probing signal (under some assumption) can connect to probability measure 
    \item See that weighting of norms connects to a way in which are probing sampling our domain 
    \item Exploration aspect in a different session: On vs off policy (theory vs improving policy)
    \item We had one example of off-policy divergence. Counter example: We have to be careful 
    \item Concluding part of this section was: what kind of divergence can arise? Deadly triad: being able to explain difficulty 
    \item Switch to other sources 
    \item MC policy evaluation aspects. Instead of TD view (1 step), I am now observing estimates of cost 
    \item We need to approximate a function, where we have some function values. At first, we solved some linear equations system 
    \item MC view gives a different way of solving this (more efficiently)
    \item Importance sampling view: Might put more importance on total cost, which are large! 
    \item We don't really want to solve equations \(\implies\) gradient methods in least square approximation problems 
    \item This gives direct policy evaluation (as before we had TD view for example). Batch versions 
    \item Rewriting gives TD error, saw this is TD(1)
    \item Used this as motivation for multistep methods 
    \item T(\(\lambda\)), l+1 step version of bellman operator to reintroduce TD\((\lambda)\)
    \item First view into stochastics: bias variance Tradeoff, \(\lambda\) is our tuing parameter, TD(1) might not be the best
    \item Final section on policy gradient methods, before we approximated Q or value function and indirectly got policy. Now we approximate policy directly
    \item What if we parametrize our policy? Stochastic policy (softmax example)
    \item Policy gradient theorems, seeing in different ways: total cost can be evaluated reasonable 
    \item Theorem to compute score?
    \item REINFORCE approach first in finite horizon, then infinite horizon, policy gradient results to show convergence results 
    \item Mini-batch REINFORCE, to get unbiased estimator, almost sure convergence for L-smooth functions (connecting back to earlier)
    \item View lemmas, assumption (QSA1, etc.)
    \item Gradient of total cost goes to zero \(\to \) function goes to critical point
    \item Also for stochastic gradient stuff we can do convergence theory 
    \item Last two lectures: Actor critic: combining policy and Q function approximation and use both 
    \item One can see AC as policy iteration, we came back to policy iteration anyway.
\end{itemize}

Overall: reasonable overview of RL both theory and praxis, still a lot missing! Nobody does linear function approximation, but rather NNs.
Hope: guide algorithm by theory.

\markeol{24}
